% Proposed B.146. Owned source; no private PDF/PNG requested.
\subsection*{All cyclic leaf relations and the exact ambient restriction}

\statement{B.146}{Proved}{The complete cyclic mixed-square and four-leaf filling is a proper internal amalgam, with an explicit all-word ambient restriction kernel}
\label{stmt:B146}
In the actual B.141 unfilled group $V=R_L*_KR_R$, write
\[
 K=F(a,B,d,e,f),\quad
 (V_1,V_2,V_3,V_4)=(a,d,e,f^{-1}Bf),\quad
 q=\eta(afa^{-1})^{-1}.
\]
Let $\mathfrak d:K\hookrightarrow R_R$ denote the complete marked
map of B.141, and put $\gamma=\mathfrak d(f)$,
$x_i=\mathfrak d(V_i)$. With $N_R=\ker(r_R:R_R\to K)$, the
actual cyclic subgroup and its complete five-relation filling are
\begin{equation}\label{eq:B146-carrier}
 \mathscr D=\langle k\rangle*B_q,\qquad
 B_q=N_R\rtimes\langle q\rangle,
\end{equation}
\begin{equation}\label{eq:B146-five}
 R_{\rm cyc}=\langle\!\langle P_0,P_1,P_2,P_3,P_4
                    \rangle\!\rangle_{\mathscr D},\quad
 P_0=[(\gamma k)^2,k],\quad P_i=[k,x_i].
\end{equation}
There is an exact proper ordinary amalgam
\begin{equation}\label{eq:B146-amalgam}
 \mathscr D/R_{\rm cyc}
       =B_q*_{F_*}E_{\rm new},\qquad
 F_*=\mathfrak d(K)=F(\gamma,x_1,x_2,x_3,x_4),
\end{equation}
where $E_{\rm new}$ is specified in the proof. Both vertex groups
embed internally. The full initial filling still uses every $K$
translate and the actual $\theta$-family; its restriction on this
carrier is exactly \eqref{eq:B146-ambient-kernel}. Neither internal
properness nor the cyclic normalization makes that kernel trivial.

\begin{proof}
\emph{B.146.1. The actual embedded carrier and its free edge.}
Use the actual common marking $p:R^0\to R_R$ with $p|_K=\eta$.
The precise map is
\begin{equation}\label{eq:B146-d-map}
 \mathfrak d(u)=p\bigl(aC(a^{-1}\eta(u)a)Ca^{-1}\bigr).
\end{equation}
It is injective and retains the whole conjugator, as proved in
B.141. Its retraction values are
$r_R(\gamma)=q$ and $r_R(x_i)=1$; thus $F_*\leq B_q$.
The tuple $(f,V_1,V_2,V_3,V_4)$ is a free basis of $K$, so the
five displayed edge generators form an actual free basis of $F_*$.

B.145 gives the actual free pair $U=\langle k,q\rangle\leq R_L$.
Its kernel over $\langle q\rangle$ has free basis
$q^nkq^{-n}$, $n\in\mathbb Z$, and lies in $N_L$.
Consequently $U\cap K=\langle q\rangle$. By its definition,
$B_q\cap K=\langle q\rangle$. Ordinary amalgam normal form
therefore embeds
\[
 U*_{\langle q\rangle}B_q
       =\langle k\rangle*B_q
\]
in the actual $V$. Indeed every reduced alternating syllable
outside the indicated edge stays outside $K$ in its respective
ambient vertex group. This proves \eqref{eq:B146-carrier}; it
retains the entire $N_R$, not a selected finite family of words.
It also agrees with the ordinary retraction in B.145.

\emph{B.146.2. The complete internal five-relation filling.}
Define
\begin{equation}\label{eq:B146-E-new}
 E_{\rm new}=\langle\gamma,x_1,x_2,x_3,x_4,k_{\rm new}
 \mid[(\gamma k_{\rm new})^2,k_{\rm new}],
       [k_{\rm new},x_i]\ (1\leq i\leq4)\rangle.
\end{equation}
The homomorphism $k_{\rm new}\mapsto1$ and the identity on the
other five letters retracts this group onto $F_*$. Hence its
edge map is injective. The substitution
$h_{\rm new}=\gamma k_{\rm new}$ also gives the exact presentation
\[
 \langle h_{\rm new},k_{\rm new}\mid
                     [h_{\rm new}^{,2},k_{\rm new}]\rangle
 *_{\langle k_{\rm new}\rangle}
 \bigl(\langle k_{\rm new}\rangle\times F(x_1,x_2,x_3,x_4)\bigr).
\]
Equality \eqref{eq:B146-amalgam} now follows directly from the
presentations of $\mathscr D=\langle k\rangle*B_q$ and
\eqref{eq:B146-E-new}, identifying $k$ with $k_{\rm new}$.
Both specified maps on $F_*$ are injective, so ordinary reduced
normal form embeds both vertex groups. There is also a retraction
to $B_q$ killing $k_{\rm new}$. Since $q\in\mathscr D$, its
normal subgroup in \eqref{eq:B146-five} contains the whole
$q^{\mathbb Z}$ orbit of all five relations. This is a complete
internal assertion, with no truncation in the exponent.

\emph{B.146.3. All four leaf orbits in the two-generator square model.}
Write the actual $\gamma=c_0q$, with $c_0\in N_R$, and let
\[
 b=\operatorname{Ad}(\gamma)|_{N_R}
       =\operatorname{Inn}(c_0)\operatorname{Ad}(q)|_{N_R}.
\]
Changing the stable letter from $q$ to this actual $\gamma$ does
not remove the correction. Put $z_n=\gamma^nk\gamma^{-n}$.
The square relation is equivalent, by cancellation of its final
$k$, to
\[
 \gamma k\gamma k=k\gamma k\gamma.
\]
Its complete integer orbit is
\begin{equation}\label{eq:B146-z-recurrence}
 z_n=z_{n+1}z_{n+2}z_{n+1}^{-1}\qquad(n\in\mathbb Z).
\end{equation}
The same two-sided inverse-presentation argument as B.145 gives
$N_R*F(z_0,z_1)$ after the square-only kernel filling. Its shift is
\begin{equation}\label{eq:B146-shift}
\begin{aligned}
 \beta|_{N_R}&=b,&\beta(z_0)&=z_1,&
 \beta(z_1)&=z_1^{-1}z_0z_1,\
 \beta^{-1}|_{N_R}&=b^{-1},&\beta^{-1}(z_1)&=z_0,&
 \beta^{-1}(z_0)&=z_0z_1z_0^{-1}.
\end{aligned}
\end{equation}
Each inverse identity follows by free-word substitution.

For every signed integer $n$, the entire word
$b_n=\gamma^nq^{-n}$ belongs to $N_R$. Thus
\begin{equation}\label{eq:B146-leaf-conjugator}
 \gamma^n[k,x_i]\gamma^{-n}
    =b_n\bigl(q^n[k,x_i]q^{-n}\bigr)b_n^{-1}.
\end{equation}
These are actual conjugators, including negative $n$, so the two
leaf normal closures in the square kernel coincide. All four
complete leaf families are therefore exactly
\begin{equation}\label{eq:B146-all-leaves}
 [z_n,b^n(x_i)]=1\qquad(n\in\mathbb Z,\ 1\leq i\leq4)
\end{equation}
in $N_R*F(z_0,z_1)$, with every $z_n$ defined by the invertible
recurrence \eqref{eq:B146-z-recurrence}. Equivalently the internal
amalgam in \eqref{eq:B146-amalgam} is
\[
 \left((N_R*F(z_0,z_1))/
       \langle\!\langle\eqref{eq:B146-all-leaves}\rangle\!\rangle\right)
       \rtimes_\beta\langle\gamma\rangle.
\]
The relation subgroup is $\beta$-invariant by its whole integer
indexing. This proves both presentation directions and retains
every nonlinear right-kernel word. It does not assert that the
leaf words are coordinate graph generators.

\emph{B.146.4. The full retained action, both mirror families and the remaining kernel.}
For each actual $g\in K$, conjugate the entire construction in
$V$. This gives $\mathscr D_g=g\mathscr Dg^{-1}$,
$B_g=N_R\rtimes\langle gqg^{-1}\rangle$, the specified edge
$gF_*g^{-1}$, and the corresponding internally proper quotient.
Apply the actual involution $\theta$ to obtain the opposite
family. This is one all-word construction, not a list of tested
translates. The full initial normal subgroup of B.141 is exactly
\begin{equation}\label{eq:B146-full-normal}
 N_*=\langle\!\langle gR_{\rm cyc}g^{-1},\
          \theta(gR_{\rm cyc}g^{-1}):g\in K
                       \rangle\!\rangle^V.
\end{equation}
Both squares and all eight leaf types are included. The exact
remaining restriction on the particular cyclic carrier is
\begin{equation}\label{eq:B146-ambient-kernel}
 \ker(\mathscr D/R_{\rm cyc}\longrightarrow V/N_*)
          =(\mathscr D\cap N_*)/R_{\rm cyc}.
\end{equation}
The proper internal model gives no vanishing of this intersection.
Other retained directions and the opposite family can impose
additional relations; pairwise subgroup intersections do not
replace their entire ambient normal closure. B.145's specified
retraction is discontinuous for the ambient full topology, so it
cannot promote this internal result to a compact filling theorem.
No completion/intersection exchange is used here.

The complete prior proof is retained at
\path{work_DE/next_marginal_source/cyclic_leaf_filling_continuation29.md},
with its independently bound full review; the exact square input
is the B.145 certificate. Every assertion above has its own
presentation or normal-form proof and does not rely on a research
file in place of a proof.
\end{proof}
